% !TeX root = main.tex
\lecture{5}{Thu 05 Feb 2026 15:00}{Mechanics V: Rigid Pendula}

Reminder, angular momentum:
\begin{itemize}
    \item A vector, perpendicular to $\vec{r}, \vec{p}$, given by $\vec{L} = \vec{r} \wedge \vec{p}$.
    \item Conserved under central forces (i.e. forces pointing radially). 
    \item Must be defined about some origin, hence is defined about a point in space.
\end{itemize}

Reminder, torque:
\begin{itemize}
    \item A vector, perpendicular to $\vec{r}, \vec{F}$. given by $\vec{\tau} = \vec{r} \wedge \vec{F}$.
    \item Also defined about a point.
    \item For a fixed force, the torque grows with distance away from the fulcrum.
    \item Can also be approached using the magnitude $\tau = I \ddot{\theta} = I\dot{\omega}$
\end{itemize}

\section{Example: Small Oscillations of a Simple Pendulum}
Consider a mass at a frictionless pivot (the origin). The bob has mass $m$ and makes angle $\theta$ with the horizontal. The rod connecting them is light.

Considering the torque about $O$:
\[
    \vec{\tau} = \vec{r} \wedge \vec{F} = \det \begin{bmatrix}
        \hat{\imath} & \hat{\jmath} & \hat{k} \\
        l \sin \theta & -l \cos \theta & 0 \\
        0 & -mg & 0
    \end{bmatrix} = \hat{\vec{k}} \left(-mg l \sin \theta\right)
\]

We also have that:
\[
    \tau = I \ddot{\theta}
\]

Hence:
\[
    I \ddot{\theta} = -mg l \sin \theta \tag{1}
\]

For a point mass, the moment of inertia at distance $R = l$ is $I = ml^2$. As we have small oscillations, $\sin \theta \approx \theta$. Hence:
\[
    ml^2 \ddot{\theta} \approx -mgl \theta
\]
\[
    \ddot{\theta} = - \frac{g \theta}{l}
\]

This gives us simple harmonic motion with:
\[
    \omega^2 = \frac{g}{l}
\]
where:
\[
    \theta = \theta_0 \cos \omega t
\]

Given that simply harmonic motion is defined by:
\[
    \ddot{x} + \omega^2 x = 0
\]

Hence $T = \frac{2 \pi}{\omega} = 2 \pi \sqrt{\frac{l}{g}}$.

\section{Rigid Pendulum}
Consider a pendulum of any arbitrary shape and mass distribution where the pivot point is the top edge of the object.

Suppose the vector from the pivot point to the centre of mass (with downward force $mg$) has length $l$ and is at angle from the downward vertical $\theta$.

We still have:
\[
    \vec{\tau} = \vec{r} \wedge \vec{F} = \vec{l} \times mg = -lmg sin \theta
\]
\[
    \tau = I_A \ddot{\theta}
\]

Hence:
\[
    I_A \ddot{\theta} = -lmg \sin \theta \approx - l m g \theta
\]

Hence again:
\[
    \ddot{\theta} = \frac{-mgl \theta}{I_A}
\]

This still gives us simple harmonic motion, but our $\omega^2$ is now $\omega^2 = mgl/I_A$. If we know the moment of inertia about the CoM, we can use $I_0$ with the parallel axes theorem hence:
\[
    \omega^2 = \frac{mgl}{I_0 + ml^2}
\]

We can also invoke the radius of gyration to make it mass independent and depend only on the shape (which is encoded into $K$):
\[
    \omega^2 = \frac{mgl}{mK^2 + ml^2}
\]
\[
    = \frac{gl}{K^2 + l^2}
\]

The time period is again:
\[
    T = \frac{2 \pi}{\omega} = 2 \pi \sqrt{\frac{K^2+l^2}{lg}}
\]

Lets check that for the simple pendulum where the shape is a point mass, $I_0 = 0 \implies k = 0$.

Hence applying the general formula to a point mass case gives:
\[
    T = 2 \pi \sqrt{l^2} / gl = 2 \pi \sqrt{\frac{l}{g}}
\]
Which is as-expected!

\section{Example: Solid Uniform Rod}
Suppose our object is a solid uniform rod of mass $M$ oscillating with angle $\theta$ and of length $d$. The CoM lies at $O$ at position $d/2$ on the rod relative to the axis $A$ at the pivot. Does this oscillate faster or slower than the point mass example?

It can be shown that the moment of inertia about the CoM is:
\[
    I_0 = Md^2 / 12
\]

Using the parallel axes theorem:
\[
    I_A = I_0 + \left(\frac{d}{2}\right)^2 M = \frac{Md^2}{12} + \frac{Md^2}{4} = \frac{Md^2}{3}
\]

We have shown that:
\[
\omega = \sqrt{\frac{mgl}{I_A}}
\]
Where (here) $l = d/2$.
\[
    \omega = \sqrt{\frac{3mgd}{2 md^2}} = \frac{3 g}{2d}
\]
This gives:
\[
    T = 2 \pi \sqrt{\frac{2}{3} \frac{d}{g}}
\]
Which is a slower oscillation compared to a simple pendulum.

\section{Massive Pulleys}
Consider a pulley system with a massless string pulled over a uniform pulley (at a fixed support) with radius $R$. The LHS dangling mass has mass $m_1$ and the RHS has a smaller mass $m_2$. The LHS has tension $T_1$ and RHS tension $T_2$. 

\textbf{Equations of Motion for the Masses}

Using $F = ma$ on the two masses gives:
\[
    m_1 g - T_1 = m_1 a \tag{1}
\]
\[
    T_2 - m_2 g = m_2 a \tag{2}
\]

\textbf{Case 1: Massless Pulley}

If the pulley is massless, the two tensions are equal $T_1 = T_2 = T$
\[
    m_1 g - T = m_1 a \tag{1b}
\]
\[
    T m_2 g = m_2 a \tag{2b}
\]

$1b + 2b$ gives:
\[
    m_1 g - m_2g = (m_1 + m_2)a
\]
\[
    \implies a = \left(\frac{m_1 - m_2}{m_1 + m_2}\right)g
\]

\textbf{Case 2: Massive Pulley}
The pulley now has mass $M$. In order to get the pulley moving, there needs to be some kind of net torque on the pulley. This means that the two tensions $T_1$ and $T_2$ can no longer be assumed to be equal.
\[
    \tau = I \ddot{\theta} \equiv I \alpha
\]
\[
    \vec{\tau} = \vec{r} \wedge \vec{F} = \vec{R} \wedge (T_2 - T_1)
\]
Hence:
\[
    I \alpha = R(T_2 - T_1) \tag{3}
\]
And as its a uniform disk:
\[
    I = \frac{MR^2}{2}
\]
And converting from angular formulae to linear formulae:
\[
    v = R \omega \implies a = R \alpha
\]
Inserting these into (3):
\[
    \frac{MR^2}{2} \frac{a}{R} = R(T_2 - T_1)
\]
\[
    \implies \frac{Ma}{2} = T_2 - T_1 \tag{4}
\]
And inserting our general equations of motion:
\[
    T_2 = m_2 g - m_2 a, \qquad T_1 = m_1 g + m_1 a
\]
\[
    \frac{Ma}{2} = (m_2 g - m_2 a) - (m_1 g + m_1 a)
\]
\[
    \frac{Ma}{2} = (m_2 - m_1)g - (m_1 + m_2)a
\]
\[
    \frac{Ma}{2} + (m_1 + m_2)a = (m_2 - m_1)g
\]
\[
    a\left(\frac{M}{2} + m_1 + m_2\right) = (m_2 - m_1)g
\]
\[
    \boxed{a = \frac{(m_2 - m_1)g}{\frac{M}{2} + m_1 + m_2}}
\]



